\documentclass[11pt,a4paper,notitlepage]{report}
% Shared preamble for the raw chapter-4 write-ups under docs/report/insights/.
%
% Each experiment_N/experiment_N.tex is a standalone, compilable LaTeX report
% written BEFORE anything goes into final_report/chapters/chapter_4_*.tex.
% It uses the final report's fonts, units and macros so that passages, tables
% and figures lift into the chapter without reformatting.
%
%   \documentclass[11pt,a4paper]{report}
%   % Shared preamble for the raw chapter-4 write-ups under docs/report/insights/.
%
% Each experiment_N/experiment_N.tex is a standalone, compilable LaTeX report
% written BEFORE anything goes into final_report/chapters/chapter_4_*.tex.
% It uses the final report's fonts, units and macros so that passages, tables
% and figures lift into the chapter without reformatting.
%
%   \documentclass[11pt,a4paper]{report}
%   % Shared preamble for the raw chapter-4 write-ups under docs/report/insights/.
%
% Each experiment_N/experiment_N.tex is a standalone, compilable LaTeX report
% written BEFORE anything goes into final_report/chapters/chapter_4_*.tex.
% It uses the final report's fonts, units and macros so that passages, tables
% and figures lift into the chapter without reformatting.
%
%   \documentclass[11pt,a4paper]{report}
%   \input{../common/preamble}
%   \begin{document} ...
%
% Figures and tables are NOT copied: \graphicspath and \reporttable resolve into
% final_report/figures/ and final_report/tables/, so the raw write-up always
% shows the same artefacts as the chapter.  Compile with:  latexmk -pdf <file>

\usepackage[T1]{fontenc}
\usepackage[utf8]{inputenc}
\usepackage{mathptmx}
\usepackage[scaled=0.92]{helvet}
\usepackage{courier}
\usepackage[a4paper,left=30mm,right=20mm,top=25mm,bottom=25mm]{geometry}
\usepackage{microtype}
\usepackage{booktabs}
\usepackage{tabularx}
\usepackage{longtable}
\usepackage{array}
\usepackage{multirow}
\usepackage{amsmath}
\usepackage{amssymb}
\usepackage{ifthen}
\usepackage{siunitx}
\usepackage{graphicx}
\usepackage{placeins}
\usepackage{caption}
\usepackage{enumitem}
\usepackage[table]{xcolor}
\usepackage{rotating}
\PassOptionsToPackage{hyphens}{url}\usepackage{url}
\usepackage[hidelinks]{hyperref}
\usepackage[nameinlink,noabbrev]{cleveref}

\graphicspath{{../../final_report/}{../../final_report/figures/}}
% \reporttable{name} inputs final_report/tables/name.tex
\newcommand{\reporttable}[1]{\input{../../final_report/tables/#1}}

\newcommand{\repo}[1]{\texttt{\detokenize{#1}}}
\sisetup{detect-all,group-separator={,},range-phrase=--,range-units=single}
\DeclareSIUnit{\pixel}{px}
\DeclareSIUnit{\mebi}{Mi}
\captionsetup{font=small,labelfont=bf}
\setlist{nosep}
\setlength{\parskip}{4pt}
\setlength{\LTcapwidth}{\textwidth}

% ---- result-slot status (RENDERED here, unlike in the final report)
\definecolor{stcompleted}{HTML}{A5D6A7}
\definecolor{stpartial}{HTML}{FFF2CC}
\definecolor{stuntested}{HTML}{FFD8A8}
\definecolor{stblocked}{HTML}{FFCDD2}
\definecolor{stfailed}{HTML}{FFCDD2}
\definecolor{stplanned}{HTML}{E0E0E0}
\newcommand{\slotstatus}[2]{%
  \par\noindent\fcolorbox{black!40}{st#1}{\parbox{\dimexpr\textwidth-2\fboxsep-2\fboxrule}{%
    \textbf{Status: #1.}\ #2}}\par\medskip}
\newcommand{\status}[1]{\textbf{#1}}

% ---- queries: RENDERED in red here so a reader sees every open concern in place
\definecolor{flagred}{HTML}{B00000}
\newcommand{\flag}[1]{\textcolor{flagred}{#1}}
\newcommand{\moved}[1]{#1}
\newenvironment{queried}[1][Questionable]%
  {\par\addvspace{0.6ex}\begingroup\color{flagred}\noindent
   \textbf{[#1]}\ \ignorespaces}%
  {\par\endgroup\addvspace{0.6ex}}

% ---- per-block performance marks (section 4.5 diagrams)
\definecolor{perfgood}{HTML}{A5D6A7}
\definecolor{perfpoor}{HTML}{FFCDD2}
\definecolor{perfnoeffect}{HTML}{D9EAF7}
\definecolor{perfnone}{HTML}{FFF2CC}
\newcommand{\perfmark}[1]{%
  \ifthenelse{\equal{#1}{good}}{\colorbox{perfgood}{\scriptsize good}}{%
  \ifthenelse{\equal{#1}{poor}}{\colorbox{perfpoor}{\scriptsize poor}}{%
  \ifthenelse{\equal{#1}{noeffect}}{\colorbox{perfnoeffect}{\scriptsize no effect}}{%
  \colorbox{perfnone}{\scriptsize no claim}}}}}

%   \begin{document} ...
%
% Figures and tables are NOT copied: \graphicspath and \reporttable resolve into
% final_report/figures/ and final_report/tables/, so the raw write-up always
% shows the same artefacts as the chapter.  Compile with:  latexmk -pdf <file>

\usepackage[T1]{fontenc}
\usepackage[utf8]{inputenc}
\usepackage{mathptmx}
\usepackage[scaled=0.92]{helvet}
\usepackage{courier}
\usepackage[a4paper,left=30mm,right=20mm,top=25mm,bottom=25mm]{geometry}
\usepackage{microtype}
\usepackage{booktabs}
\usepackage{tabularx}
\usepackage{longtable}
\usepackage{array}
\usepackage{multirow}
\usepackage{amsmath}
\usepackage{amssymb}
\usepackage{ifthen}
\usepackage{siunitx}
\usepackage{graphicx}
\usepackage{placeins}
\usepackage{caption}
\usepackage{enumitem}
\usepackage[table]{xcolor}
\usepackage{rotating}
\PassOptionsToPackage{hyphens}{url}\usepackage{url}
\usepackage[hidelinks]{hyperref}
\usepackage[nameinlink,noabbrev]{cleveref}

\graphicspath{{../../final_report/}{../../final_report/figures/}}
% \reporttable{name} inputs final_report/tables/name.tex
\newcommand{\reporttable}[1]{\input{../../final_report/tables/#1}}

\newcommand{\repo}[1]{\texttt{\detokenize{#1}}}
\sisetup{detect-all,group-separator={,},range-phrase=--,range-units=single}
\DeclareSIUnit{\pixel}{px}
\DeclareSIUnit{\mebi}{Mi}
\captionsetup{font=small,labelfont=bf}
\setlist{nosep}
\setlength{\parskip}{4pt}
\setlength{\LTcapwidth}{\textwidth}

% ---- result-slot status (RENDERED here, unlike in the final report)
\definecolor{stcompleted}{HTML}{A5D6A7}
\definecolor{stpartial}{HTML}{FFF2CC}
\definecolor{stuntested}{HTML}{FFD8A8}
\definecolor{stblocked}{HTML}{FFCDD2}
\definecolor{stfailed}{HTML}{FFCDD2}
\definecolor{stplanned}{HTML}{E0E0E0}
\newcommand{\slotstatus}[2]{%
  \par\noindent\fcolorbox{black!40}{st#1}{\parbox{\dimexpr\textwidth-2\fboxsep-2\fboxrule}{%
    \textbf{Status: #1.}\ #2}}\par\medskip}
\newcommand{\status}[1]{\textbf{#1}}

% ---- queries: RENDERED in red here so a reader sees every open concern in place
\definecolor{flagred}{HTML}{B00000}
\newcommand{\flag}[1]{\textcolor{flagred}{#1}}
\newcommand{\moved}[1]{#1}
\newenvironment{queried}[1][Questionable]%
  {\par\addvspace{0.6ex}\begingroup\color{flagred}\noindent
   \textbf{[#1]}\ \ignorespaces}%
  {\par\endgroup\addvspace{0.6ex}}

% ---- per-block performance marks (section 4.5 diagrams)
\definecolor{perfgood}{HTML}{A5D6A7}
\definecolor{perfpoor}{HTML}{FFCDD2}
\definecolor{perfnoeffect}{HTML}{D9EAF7}
\definecolor{perfnone}{HTML}{FFF2CC}
\newcommand{\perfmark}[1]{%
  \ifthenelse{\equal{#1}{good}}{\colorbox{perfgood}{\scriptsize good}}{%
  \ifthenelse{\equal{#1}{poor}}{\colorbox{perfpoor}{\scriptsize poor}}{%
  \ifthenelse{\equal{#1}{noeffect}}{\colorbox{perfnoeffect}{\scriptsize no effect}}{%
  \colorbox{perfnone}{\scriptsize no claim}}}}}

%   \begin{document} ...
%
% Figures and tables are NOT copied: \graphicspath and \reporttable resolve into
% final_report/figures/ and final_report/tables/, so the raw write-up always
% shows the same artefacts as the chapter.  Compile with:  latexmk -pdf <file>

\usepackage[T1]{fontenc}
\usepackage[utf8]{inputenc}
\usepackage{mathptmx}
\usepackage[scaled=0.92]{helvet}
\usepackage{courier}
\usepackage[a4paper,left=30mm,right=20mm,top=25mm,bottom=25mm]{geometry}
\usepackage{microtype}
\usepackage{booktabs}
\usepackage{tabularx}
\usepackage{longtable}
\usepackage{array}
\usepackage{multirow}
\usepackage{amsmath}
\usepackage{amssymb}
\usepackage{ifthen}
\usepackage{siunitx}
\usepackage{graphicx}
\usepackage{placeins}
\usepackage{caption}
\usepackage{enumitem}
\usepackage[table]{xcolor}
\usepackage{rotating}
\PassOptionsToPackage{hyphens}{url}\usepackage{url}
\usepackage[hidelinks]{hyperref}
\usepackage[nameinlink,noabbrev]{cleveref}

\graphicspath{{../../final_report/}{../../final_report/figures/}}
% \reporttable{name} inputs final_report/tables/name.tex
\newcommand{\reporttable}[1]{\input{../../final_report/tables/#1}}

\newcommand{\repo}[1]{\texttt{\detokenize{#1}}}
\sisetup{detect-all,group-separator={,},range-phrase=--,range-units=single}
\DeclareSIUnit{\pixel}{px}
\DeclareSIUnit{\mebi}{Mi}
\captionsetup{font=small,labelfont=bf}
\setlist{nosep}
\setlength{\parskip}{4pt}
\setlength{\LTcapwidth}{\textwidth}

% ---- result-slot status (RENDERED here, unlike in the final report)
\definecolor{stcompleted}{HTML}{A5D6A7}
\definecolor{stpartial}{HTML}{FFF2CC}
\definecolor{stuntested}{HTML}{FFD8A8}
\definecolor{stblocked}{HTML}{FFCDD2}
\definecolor{stfailed}{HTML}{FFCDD2}
\definecolor{stplanned}{HTML}{E0E0E0}
\newcommand{\slotstatus}[2]{%
  \par\noindent\fcolorbox{black!40}{st#1}{\parbox{\dimexpr\textwidth-2\fboxsep-2\fboxrule}{%
    \textbf{Status: #1.}\ #2}}\par\medskip}
\newcommand{\status}[1]{\textbf{#1}}

% ---- queries: RENDERED in red here so a reader sees every open concern in place
\definecolor{flagred}{HTML}{B00000}
\newcommand{\flag}[1]{\textcolor{flagred}{#1}}
\newcommand{\moved}[1]{#1}
\newenvironment{queried}[1][Questionable]%
  {\par\addvspace{0.6ex}\begingroup\color{flagred}\noindent
   \textbf{[#1]}\ \ignorespaces}%
  {\par\endgroup\addvspace{0.6ex}}

% ---- per-block performance marks (section 4.5 diagrams)
\definecolor{perfgood}{HTML}{A5D6A7}
\definecolor{perfpoor}{HTML}{FFCDD2}
\definecolor{perfnoeffect}{HTML}{D9EAF7}
\definecolor{perfnone}{HTML}{FFF2CC}
\newcommand{\perfmark}[1]{%
  \ifthenelse{\equal{#1}{good}}{\colorbox{perfgood}{\scriptsize good}}{%
  \ifthenelse{\equal{#1}{poor}}{\colorbox{perfpoor}{\scriptsize poor}}{%
  \ifthenelse{\equal{#1}{noeffect}}{\colorbox{perfnoeffect}{\scriptsize no effect}}{%
  \colorbox{perfnone}{\scriptsize no claim}}}}}

\setcounter{chapter}{4}
\renewcommand{\thesection}{\arabic{section}}
\renewcommand{\thefigure}{I.\arabic{figure}}
\renewcommand{\thetable}{I.\arabic{table}}
% every table and figure of this document is generated by script/exp4_insights.py
\newcommand{\insdir}{../../../../output/compare/simulation/exp4-insights}
\newcommand{\ins}[1]{\input{\insdir/#1}}

\title{Experiment 4: Where the Map Is Wrong\\
       \large Insight analysis for Chapter 4, Section 4.4}
\author{Bolton Athit Davies}
\date{Reporting cutoff: 26 October 2026}

\begin{document}
\maketitle

\noindent\textbf{Purpose of this document.} The raw write-up of Experiment~4
(\path{insights/experiment_4/}) reports that \repo{map/map_tiny.db} places
\SI{96.6}{\percent} of its occupied cells on real structure and reproduces
\SI{89.5}{\percent} of the surfaces it observed. This document asks what the
remaining \SI{3.4}{\percent} and \SI{10.5}{\percent} are, and whether the map's
quality follows anything a reader would propose: range from the route, how often
a surface was observed, the angle it was seen at, the loop-closure correction,
the kind of object. It also opens the four unregistered databases that sit
beside the scored one. Every number is written by \path{script/exp4_insights.py}
into \path{output/compare/simulation/exp4-insights/}; the tables below are
input from there unchanged. One warehouse and two maps support descriptive
patterns, not statistics, and the document says so where it matters. Added 26
October 2026: the same classification applied to the two realistic-warehouse
maps of the raw write-up (Insight~6), written by the script's \repo{--realistic}
entry into files named \repo{*_realistic.*} beside the tiny ones, which it does
not rewrite.

\section{Questions and method}

\begin{enumerate}[label=Q\arabic*.]
\item What are the wrong occupied cells: surfaces drawn too thick, or obstacles
      where there is nothing?
\item What are the missed surfaces, and does the miss follow range, the number
      of observing nodes, the viewing angle, or the object?
\item Did the loop closures change the map's accuracy?
\item On an independent route, where does the map block the robot, and why?
\item What are the four unregistered databases, and which recording built the
      scored map?
\item Do the two realistic-warehouse maps, four times the area and built
      without loop closure, fail by the same mechanism, and what blocks the
      independent probe on them?
\end{enumerate}

\paragraph{Raster.} The conventions are those of \path{script/validate_map.py}:
the exported occupancy grid is placed in the world by the map pose that
\path{script/fit_map_pose.py} measures from the grid alone, true structure is
the world file's obstacle footprints filled plus the wall ring at \SI{0.05}{\metre}
per cell, true \emph{surfaces} are the one-cell boundary of that, and a cell is
correct within \SI{0.25}{\metre}. Occupied cells beyond the tolerance are split
at \SI{0.50}{\metre}: nearer is a surface drawn too thick, further is a phantom.
Recall is scored on the surfaces the map observed, as the chapter does.

\paragraph{Per-node observation.} RTAB-Map stores, for every graph node, the
local occupancy cells it contributed (obstacle, ground and empty, in the node
frame). Moving them through the optimised node poses gives, for every world
cell, how many nodes observed it and which node observed it first. The
incidence angle of a surface is the angle between its outward normal and the
line of sight from the nearest graph node: \SI{0}{\degree} faced head-on,
\SI{90}{\degree} seen edge-on along its own plane.

\paragraph{Second map.} \repo{map/map_02.db}, one of the unregistered
databases, fits the tiny world as well as the scored map does (Q5) and was
built on a different route. It is scored by the same method and used as the
independent drivability probe that the lost \repo{dataset_static_tiny_repeat}
reference was meant to provide.

\FloatBarrier
\section{Insight 1: the wrong cells are thick surfaces, not phantom obstacles}

Of the \SI{3.4}{\percent} of occupied cells that miss the tolerance, nine in
ten lie within \SI{0.50}{\metre} of a real surface, and the rest form 22
clusters of at most ten cells. \Cref{tab:insight4-classes} gives the split and
\cref{fig:cells-tiny} the places.

\ins{cell_classes.tex}

\begin{figure}[p]
  \centering
  \includegraphics[width=\textwidth]{\insdir/fig_cells_map_tiny_light.png}
  \caption{\repo{map_tiny.db} against the tiny world: correct occupied cells
  (green), surfaces drawn too thick (orange), phantom cells (red, the largest
  clusters labelled with their cell count), true surfaces the map observed but
  did not reproduce (blue), the recorded route and the optimised graph nodes.}
  \label{fig:cells-tiny}
\end{figure}

The thickened cells sit on the wall ring and on shelf fronts, where the stereo
depth of a planar surface spreads over two or three cells. The 43 phantom cells
together cover about \SI{0.1}{\square\metre}; the largest cluster, ten cells at
the north wall near the spawn, is \SI{0.9}{\metre} from the route and was
observed by one node. None of the phantoms lies in an aisle the recorded route
uses, which is why the build route is fully drivable on this map.

\FloatBarrier
\section{Insight 2: surfaces are missed when they are seen edge-on and only
         once, not when they are far}

Recall of the observed surfaces falls from \SI{95}{\percent} for surfaces faced
head-on to \SI{76}{\percent} for surfaces seen at more than \SI{60}{\degree}
incidence, and from \SI{97}{\percent} for cells observed by three or more
nodes to \SI{78}{\percent} for cells observed by one; it does not fall with
distance from the route within the \SI{6}{\metre} the map reaches.
\Cref{tab:insight4-angle,tab:insight4-views,tab:insight4-dist} and
\cref{fig:bins} carry the three contrasts.

\ins{bins_angle.tex}
\ins{bins_views.tex}
\ins{bins_distance.tex}

\begin{figure}[htbp]
  \centering
  \includegraphics[width=\textwidth]{\insdir/fig_bins_light.png}
  \caption{Precision of the occupied cells and recall of the observed surfaces
  against distance from the recorded route (left) and against the number of
  graph nodes that observed the cell (right), for both tiny-warehouse maps.}
  \label{fig:bins}
\end{figure}

The two mechanisms are the same surfaces seen from the same place. The missed
cells in \cref{fig:cells-tiny} are the faces of neighbouring shelves across the
narrow gaps on the east side, which the route looks into from the aisle but
never enters: each face is seen once, from a node in the aisle, at grazing
incidence. Precision tells the same story from the other side: cells observed
by one node are \SI{91}{\percent} correct, cells observed by three or more
\SI{99}{\percent}. Stereo range error, the mechanism a reader would propose
first, does not show: precision stays at \SIrange{96}{98}{\percent} out to
\SI{4}{\metre} and recall has no trend with distance.

\Cref{tab:insight4-objects} lists the same effect by object. The objects with
the lowest recall were each observed by a single node, while the one large
shelf observed by four nodes is reproduced at \SI{99}{\percent}. Which objects
these are follows from the route, not from the objects: the figure-eight and
the one straight along the east wall decide what is faced and what is grazed.

\ins{missed_by_obstacle.tex}

\FloatBarrier
\section{Insight 3: the loop closures corrected the graph, and the map did not
         need it}

The five closure links corrected the odometry by at most \SI{0.26}{\metre} in
the plane and \SI{0.52}{\metre} in height, and the occupied cells first
observed after the closing revisit are as accurate as those observed before
it. \Cref{tab:insight4-closures} and \cref{fig:graph} carry the numbers.

\ins{closures.tex}

\begin{figure}[htbp]
  \centering
  \includegraphics[width=\textwidth]{\insdir/fig_graph_light.png}
  \caption{\repo{map_tiny.db}: odometry poses and optimised poses of the 29
  nodes in the optimised graph, with the loop-closure links (left), and the
  planar and vertical correction each node received (right).}
  \label{fig:graph}
\end{figure}

The correction is real but small, and it is mostly vertical: VINS-Fusion's
odometry climbed \SI{0.5}{\metre} over the \SI{70}{\metre} route, and the
closures at nodes 22 and 23, which revisit nodes 9 and 10, pulled the second
half of the graph back onto the floor. In the plane the correction stays
under the \SI{0.25}{\metre} tolerance for every node, so the cells placed from
corrected and from uncorrected poses score alike. On a route of this length the
map's accuracy was set by the odometry, with the closures as insurance that
was not drawn on.

\FloatBarrier
\section{Insight 4: on an independent route the second map blocks the robot at
         its phantom cells}

Under the true footprint along the recorded route, \repo{map_02.db} blocks
\SI{7.6}{\percent} of the poses and leaves \SI{5.5}{\percent} unmapped, and
every blocked stretch coincides with one of its phantom clusters.
\Cref{tab:insight4-driv,tab:insight4-blocked} and \cref{fig:cells-02} carry
the evidence.

\ins{drivability.tex}
\ins{blocked_segments.tex}

\begin{figure}[p]
  \centering
  \includegraphics[width=\textwidth]{\insdir/fig_cells_map_02_light.png}
  \caption{\repo{map_02.db}, unregistered, against the tiny world, placed by
  its own fit; same conventions as \cref{fig:cells-tiny}. The recorded route
  drawn is that of \repo{dataset_map_tiny}, which this map did not see.}
  \label{fig:cells-02}
\end{figure}

This is the measurement the raw write-up could not make. The build route is
drivable on the map it built by construction, since the camera cleared the
cells the robot passed through; an independent route meets the cells the
build route only looked at. On the second map the three blocked stretches sit
on its phantom clusters of 3, 16 and 3 cells, in the figure-eight and at the
start of the east straight; the route centre comes within \SI{0.0}{\metre} to
\SI{0.2}{\metre} of them. Phantoms that cover \SI{0.1}{\square\metre} of a
\SI{290}{\square\metre} floor are negligible for geometry and decisive for
drivability when they fall in an aisle.

\begin{queried}[Placement of the second map]
\repo{map_02.db} is placed by its own grid fit (agreement 0.385, margin
\SI{54}{\percent} over the best rival yaw), not by a known spawn. A placement
error of a few centimetres moves the thickened wall cells but not the phantom
clusters, so the blocked stretches stand; the per-cell precision of this map
carries that uncertainty.
\end{queried}

\FloatBarrier
\section{Insight 5: the scored map was built from the recorded tiny bag, and
         the large-warehouse attempt left a database}

The nodes of \repo{map_tiny.db} coincide with the ground truth of
\repo{dataset_map_tiny} at the same stamps to a median \SI{0.05}{\metre}, so
the map was built from that recording and the write-up's self-check is a
true self-check; and \repo{map_03.db} is a \SI{76}{\metre} run whose poses
climb to \SI{9.8}{\metre} above the floor, the large-warehouse attempt that
the chapter reports as having left no artefact. \Cref{tab:insight4-dbs}
describes all five databases.

\ins{databases.tex}

\repo{map_02.db} fits the tiny world as well as the scored map does and its
nodes do not follow the recorded route (median \SI{4.1}{\metre} from the
ground truth at the same stamps): it is a second map of the same world on a
different route, made on 21 September, and the only candidate for the lost
\repo{dataset_static_tiny_repeat} session. \repo{map_03.db} cannot be placed
in any large world (agreement 0.19, margin \SI{7}{\percent}); its 18 nodes run
\SI{37}{\metre} along one aisle and leave the floor, which is tracking failure,
not mapping failure. \repo{map_01.db} and \repo{map_02_vins.db} are two-node
starts.

\begin{queried}[Provenance]
The identifications of \repo{map_02.db} and \repo{map_03.db} rest on their
contents alone; no record names them. The raw write-up's statement that the
scored map's node stamps match no recorded dataset is contradicted by the
stamp-and-pose match above and should be withdrawn.
\end{queried}

\FloatBarrier
\section{Insight 6: on the realistic maps the shelves are thickened, the walls
         are displaced, and the cells that block the probe were laid during
         initialisation}

Of the occupied cells of the two realistic maps, 13.4 and \SI{9.7}{\percent}
are surfaces drawn too thick and 5.2 and \SI{5.3}{\percent} are phantoms,
thirteen times the tiny map's phantom share; every probe pose that the maps
block lies within \SI{0.6}{\metre} of a phantom cell, and 96 and
\SI{63}{\percent} of those blocking cells were first observed by one of the
first five graph nodes. \Cref{tab:insight4-classes-realistic} gives the split,
\cref{fig:cells-realistic,fig:cells-realistic-01} the places, and
\cref{tab:insight4-driv-realistic,tab:insight4-blocked-realistic} the blocked
stretches.

\ins{cell_classes_realistic.tex}

\begin{figure}[p]
  \centering
  \includegraphics[width=\textwidth]{\insdir/fig_cells_map_realistic_light.png}
  \caption{\repo{map_realistic.db} against the U-layout realistic world, placed
  by its fit; conventions as \cref{fig:cells-tiny}. The route drawn is the
  independent probe \repo{realistic_000}, which did not build the map; the
  graph nodes are the build route. Phantom clusters are labelled with their
  cell count.}
  \label{fig:cells-realistic}
\end{figure}

\begin{figure}[p]
  \centering
  \includegraphics[width=\textwidth]{\insdir/fig_cells_map_realistic_01_light.png}
  \caption{\repo{map_realistic_01.db} against the I-layout realistic world,
  placed by its fit; the route drawn is the probe \repo{realistic_001}.}
  \label{fig:cells-realistic-01}
\end{figure}

The thickening is the tiny map's mechanism at the shelves, and the phantoms are
a different one at the walls. In both figures the shelf rows the build route
passed are green with an orange rim, surfaces spread over two or three cells as
on the tiny map; the phantoms form a continuous red band the full length of the
east wall about \SI{0.6}{\metre} inside its true line, and a second band along
the south wall outside it, in both maps alike. The largest clusters, 108 and 70
cells on the U map and 157 and 117 on the I map, all sit on that east band, 6
to \SI{11}{\metre} from the probe route and first observed at median node 32
and 36. The walls are the structure furthest from any build node; what moved
them is not identified here.

\ins{bins_views_realistic.tex}
\ins{bins_angle_realistic.tex}
\ins{bins_distance_realistic.tex}

\begin{figure}[htbp]
  \centering
  \includegraphics[width=\textwidth]{\insdir/fig_bins_realistic_light.png}
  \caption{Precision of the occupied cells and recall of the observed surfaces
  against distance from the independent probe route (left, a proxy only: the
  probe did not build either map) and against the number of graph nodes that
  observed the cell (right), for the two realistic-warehouse maps.}
  \label{fig:bins-realistic}
\end{figure}

The view-count mechanism of Insight~2 holds and is larger. Cells observed by
one node are 66 and \SI{83}{\percent} correct, cells observed by five or more
99 and \SI{100}{\percent}, against 91 and \SI{99}{\percent} on the tiny map;
recall of the observed surfaces rises from 72 and \SI{76}{\percent} at one node
to 91 and \SI{88}{\percent} at eight or more. The incidence effect also holds,
85 and \SI{83}{\percent} head-on against 75 and \SI{78}{\percent} grazing, but
the head-on figure is ten points below the tiny map's \SI{95}{\percent}, so
viewing geometry explains the spread between surfaces and not the level they
all sit at.

\ins{drivability_realistic.tex}
\ins{blocked_realistic.tex}

The probes are blocked by the phantoms, not by the thickened shelving, and the
blocking phantoms are the oldest cells in each map. Of 134 and 213 blocked
poses, 134 and 212 lie within \SI{0.6}{\metre} of a phantom cell and none
within that distance of a thickened cell alone. The largest stretch on each map
is the first metres of the probe at the spawn, 77 and 150 poses; the 77 and 41
phantom cells under the blocked stretches have a median first-observing node
of 1, and 96 and \SI{63}{\percent} of them were first seen by node 5 or
earlier, against 6 and \SI{3}{\percent} of all phantom cells. The raw write-up
records that node~1 of each database is not at identity and that its height
settles over the first nodes, the inertial initialisation transient; the cells
laid during that transient are the ones under the robot's starting position,
which is where every probe begins. The remaining stretches, at the left shelf
column on both maps and in the one-aisle gap at $x\approx10$ on the I map, sit
on small phantom clusters laid mid-run.

\begin{queried}[Two rasters, one disagreement]
The chapter's \repo{validate_map.py} reports \SI{3.8}{\percent} blocked poses
on \repo{map_realistic}; this document's raster, whose origin is the wall ring
rather than the union of map, route and wall, reports \SI{2.1}{\percent}. The
difference is the spawn stretch alone: the chapter counts probe poses 0--183,
this raster 108--184, and the route centre there has \SI{0.00}{\metre}
clearance, so the count flips on a sub-cell offset. The other two stretches,
\repo{map_realistic_01} in full (4.4 against \SI{4.4}{\percent}), and every
clearance and unmapped figure agree to \SI{0.01}{\metre}. The chapter's figure
stands as the registered one; the spawn-stretch count carries about
\SI{1.7}{\percent} of raster sensitivity.
\end{queried}

\begin{queried}[Wall displacement, untested]
The east band is inside the wall and the south band outside it in both maps,
which a uniform scale error would not produce; drift accumulated along the
closure-free build route is the candidate, but no test separates it from a
range-dependent stereo bias at the 3 to \SI{5}{\metre} the walls were seen
from. The build-route ground truth on the removable drive would decide it.
\end{queried}

\FloatBarrier
\section{What these data cannot answer}

\begin{itemize}
\item \textbf{Whether the patterns generalise.} Two maps of one small world;
      the contrasts are descriptive. The view-count and incidence effects are
      consistent between the two maps and with how stereo mapping works, which
      is the reason to trust them, not a test.
\item \textbf{The old large-warehouse attempt.} \repo{map_03.db} exists but
      cannot be scored: no placement fits, and the pose graph left the floor.
      The two realistic-warehouse maps of Insight~6 now stand as the large
      result; whether their patterns hold under loop closure is untested,
      because they were built without it.
\item \textbf{Whether the second map's phantoms are typical.} Their cause is
      not identified; they appear where the second map's own route started
      and turned, which suggests transient depth at initialisation, untested.
\item \textbf{Distance beyond \SI{6}{\metre}.} The map reaches no further, so
      the absence of a range effect is bounded to that range.
\end{itemize}

\section{Summary}

\begin{table}[htbp]
\centering\small
\caption{Questions, answers and the strength of the evidence.}
\label{tab:insight4-summary}
\begin{tabularx}{\textwidth}{lXXl}
\toprule
 & Answer & Evidence & Confidence \\
\midrule
Q1 wrong cells & Thick surfaces, not phantoms: 348 thickened against 43 phantom cells in 22 small clusters & \cref{tab:insight4-classes}, \cref{fig:cells-tiny} & high \\
Q2 missed surfaces & Seen edge-on and once: recall 95\,\% head-on to 76\,\% grazing, 97\,\% at three nodes to 78\,\% at one; no range effect within 6\,m & \cref{tab:insight4-angle,tab:insight4-views,tab:insight4-dist} & medium, two maps \\
Q3 loop closures & Corrected 0.26\,m planar and 0.52\,m vertical; map accuracy unchanged across the revisit & \cref{tab:insight4-closures} & medium \\
Q4 independent route & Second map blocks 7.6\,\% of poses, all at its phantom clusters & \cref{tab:insight4-driv,tab:insight4-blocked} & medium, placement by fit \\
Q5 provenance & Scored map built from \repo{dataset_map_tiny}; \repo{map_02} a second tiny map; \repo{map_03} the large attempt, unscorable & \cref{tab:insight4-dbs} & high / medium \\
Q6 realistic maps & Shelves thickened as on the tiny map; walls displaced about 0.5\,m as phantom bands; probes blocked only at phantoms, 96\,\% / 63\,\% of them laid by node 5 & \cref{tab:insight4-classes-realistic,tab:insight4-blocked-realistic} & medium, one map per layout \\
\bottomrule
\end{tabularx}
\end{table}

\paragraph{Consequences for Chapter 4.} Report the map's error as surface
thickness plus a handful of phantom cells, and the recall gap as a
viewing-geometry effect of the route, not a sensor-range effect. Quote the
independent-route drivability from \repo{map_02.db} on the recorded route,
with its provenance stated, in place of the lost reference. Correct the
statement that the scored map matches no recorded dataset. Report
\repo{map_03.db} as the retained artefact of the large-warehouse attempt and
its failure as tracking loss, with the database available for a future placement
attempt. For the realistic maps, attribute the false obstacles to phantom cells
laid during initialisation and at the walls, not to thickened shelving, and
state that discarding the cells of the first few nodes, or starting the map
after the inertial transient has settled, is the one correction the data
predict would remove the spawn stretch, 57 and \SI{70}{\percent} of the blocked
poses.

\end{document}
